\documentclass[10pt,a4paper]{amsart}
\usepackage{amsmath,amssymb,amsthm,graphicx}
\usepackage{verbatim}
\newenvironment{reference}{\comment}{\endcomment}
%\newenvironment{reference}{}{}
\newenvironment{slogan}{\comment}{\endcomment}
\newenvironment{history}{\comment}{\endcomment}

% For commutative diagrams we use Xy-pic
\usepackage[all]{xy}

% We use 2cell for 2-commutative diagrams.
\xyoption{2cell}
% Theorem environments.
%
\theoremstyle{plain}
\newtheorem{theorem}[subsection]{Theorem}
\newtheorem{proposition}[subsection]{Proposition}
\newtheorem{lemma}[subsection]{Lemma}

\theoremstyle{definition}
\newtheorem{definition}[subsection]{Definition}
\newtheorem{example}[subsection]{Example}
\newtheorem{exercise}[subsection]{Exercise}
\newtheorem{situation}[subsection]{Situation}

\theoremstyle{remark}
\newtheorem{remark}[subsection]{Remark}
\newtheorem{remarks}[subsection]{Remarks}

\numberwithin{equation}{subsection}

% Macros
%
\def\lim{\mathop{\mathrm{lim}}\nolimits}
\def\colim{\mathop{\mathrm{colim}}\nolimits}
\def\Spec{\mathop{\mathrm{Spec}}}
\def\Hom{\mathop{\mathrm{Hom}}\nolimits}
\def\Ext{\mathop{\mathrm{Ext}}\nolimits}
\def\SheafHom{\mathop{\mathcal{H}\!\mathit{om}}\nolimits}
\def\SheafExt{\mathop{\mathcal{E}\!\mathit{xt}}\nolimits}
\def\Sch{\mathit{Sch}}
\def\Mor{\mathop{\mathrm{Mor}}\nolimits}
\def\Ob{\mathop{\mathrm{Ob}}\nolimits}
\def\Sh{\mathop{\mathit{Sh}}\nolimits}
\def\NL{\mathop{N\!L}\nolimits}
\def\CH{\mathop{\mathrm{CH}}\nolimits}
\def\proetale{{pro\text{-}\acute{e}tale}}
\def\etale{{\acute{e}tale}}
\def\QCoh{\mathit{QCoh}}
\def\Ker{\mathop{\mathrm{Ker}}}
\def\Im{\mathop{\mathrm{Im}}}
\def\Coker{\mathop{\mathrm{Coker}}}
\def\Coim{\mathop{\mathrm{Coim}}}

% Boxtimes
%
\DeclareMathSymbol{\boxtimes}{\mathbin}{AMSa}{"02}

%
% Macros for moduli stacks/spaces
%
\def\QCohstack{\mathcal{QC}\!\mathit{oh}}
\def\Cohstack{\mathcal{C}\!\mathit{oh}}
\def\Spacesstack{\mathcal{S}\!\mathit{paces}}
\def\Quotfunctor{\mathrm{Quot}}
\def\Hilbfunctor{\mathrm{Hilb}}
\def\Curvesstack{\mathcal{C}\!\mathit{urves}}
\def\Polarizedstack{\mathcal{P}\!\mathit{olarized}}
\def\Complexesstack{\mathcal{C}\!\mathit{omplexes}}
% \Pic is the operator that assigns to X its picard group, usage \Pic(X)
% \Picardstack_{X/B} denotes the Picard stack of X over B
% \Picardfunctor_{X/B} denotes the Picard functor of X over B
\def\Pic{\mathop{\mathrm{Pic}}\nolimits}
\def\Picardstack{\mathcal{P}\!\mathit{ic}}
\def\Picardfunctor{\mathrm{Pic}}
\def\Deformationcategory{\mathcal{D}\!\mathit{ef}}
\title{Rieffel's theorem on the bicommutant of a right ideal}
\author{The Stacks Project Authors}
\date{stacks-project:brauer.tex:local-snapshot-20261012; GNU FDL 1.2 or later}
\begin{document}
\maketitle
\noindent\textit{Editorial note.} All delivered mathematical text is the author's own, copied from the pinned Stacks Project chapter file sites-cohomology.tex; no statement or proof step was rewritten, reordered or removed.\par\medskip
\noindent\textit{Editorial note.} The proof is the author's complete argument; the author even pauses in a footnote to fix the convention for the right module structure over the bicommutant, and every step of the argument is written out in the counted span.\par\medskip
\noindent\textit{Editorial note.} Every labelled reference occurring inside the counted spans is carried in assembly.external\_references as a non-load-bearing navigation pointer (same-chapter ones at their definition site, cross-chapter ones at the invocation inside the counted span), so no load-bearing mathematics is replaced by a stub. No bibliography citation occurs inside the counted spans.\par\medskip
\noindent\textit{Editorial note.} The same-name check for this card scanned every staged Stacks chapter of this repository -- all sources/stacks-*/pinned-original-source.tex files, the card's own chapter excluded -- and the labels found are recorded in dedup.same\_name\_cross\_chapter on the card.\par\medskip
\noindent\textit{Editorial note.} The preamble is the chapter's own retained minimal preamble (sources/stacks-sites-cohomology/retained-preamble.tex), a verbatim excerpt of the native preamble: the Xy-pic setup and the native theorem declarations with the mathematical macros.\par\medskip
\noindent\textit{Allowed standard prerequisites (editorial).}\begin{itemize}
\item Endomorphism rings and bicommutants. Same chapter, the definitions of the endomorphism ring of a module and of the bicommutant, and the convention for the module structure over them. Supplies the objects A' and A'' of the proof and the side on which the bicommutant acts.
\item Two-sided ideals and the product ideal. Same chapter, the elementary theory of two-sided ideals in a ring. Supplies that the product of the ring with a nonzero right ideal is a two-sided ideal, and that a ring with no nontrivial two-sided ideal has an injective structural map determined by the identity.
\end{itemize}\medskip
\begin{lemma}
\label{lemma-rieffel}
Let $A$ be a possibly noncommutative ring with $1$ which contains no
nontrivial two-sided ideal. Let $M$ be a nonzero right ideal in $A$,
and view $M$ as a right $A$-module. Then $A$ coincides with the
bicommutant of $M$.
\end{lemma}

\begin{proof}
Let $A' = \text{End}_A(M)$, so $M$ is a left $A'$-module.
Set $A'' = \text{End}_{A'}(M)$ (the bicommutant of $M$).
We view $M$ as a right $A''$-module\footnote{This
means that given $a'' \in A''$ and $m \in M$ we have a product
$m a'' \in M$. In particular, the multiplication in $A''$
is the opposite of what you'd get if you wrote elements of $A''$
as endomorphisms acting on the left.}.
Let $R : A \to A''$ be the natural homomorphism such that
$mR(a) = ma$. Then $R$ is injective, since $R(1) = \text{id}_M$
and $A$ contains no nontrivial two-sided ideal. We claim that $R(M)$
is a right ideal in $A''$. Namely, $R(m)a'' = R(ma'')$ for $a'' \in A''$
and $m$ in $M$, because {\it left} multiplication of $M$ by any element $n$
of $M$ represents an element of $A'$, and so
$(nm)a'' = n(ma'')$ for all $n$ in $M$.
Finally, the product ideal $AM$ is a two-sided ideal, and so
$A = AM$. Thus $R(A) = R(A)R(M)$, so that $R(A)$ is a right ideal in $A''$.
But $R(A)$ contains the identity element of $A''$, and so $R(A) = A''$.
\end{proof}

\end{document}
